\section{Discussion}
\label{sec:discussion}

\subsection{The Storage vs. Computation Asymmetry on Mobile Platforms}
A central finding of our empirical investigation is the dramatic asymmetry between mobile storage bandwidth and mobile CPU inference throughput:
\begin{itemize}
    \item \textbf{Dense CPU Prefill Throughput:} On the Snapdragon 4 Gen 2 SoC, evaluating a 395-token prompt using 4 threads requires approximately $54.2$~seconds, corresponding to an effective prompt-prefill throughput of only \textbf{7.29 tokens/second}. This compute bottleneck is exacerbated by thermal throttling during sustained high-intensity GEMM operations on the Cortex-A78 cores.
    \item \textbf{Physical Flash I/O Bandwidth:} In contrast, reading a 40~MB serialized KV-cache binary from the handset's internal UFS 2.2 flash storage requires an average of \textbf{256.0~milliseconds} (and as low as 58~ms when warm).
\end{itemize}

This represents an operational asymmetry exceeding \textbf{150:1} in favor of storage restoration over online computation. While server-class deployments with high-end GPUs (e.g., NVIDIA H100) achieve prompt-prefill throughputs exceeding thousands of tokens per second—making disk I/O relatively less attractive—mobile platforms operate in a fundamentally inverted regime. On mobile CPUs, persistent flash storage serves as an ultra-fast, energy-efficient cache for expensive floating-point computations.

\subsection{Mobile Memory Pressure and LMK Resilience}
In commercial mobile operating systems, application lifecycle is governed by the Android Low Memory Killer daemon (LMKD). On budget and mid-range devices equipped with 4~GB of total physical RAM, the Android system, SurfaceFlinger compositor, and background services typically consume 1.8 to 2.2~GB.

During our physical benchmark runs, the \texttt{com.chat} process exhibited a Total Proportional Set Size (PSS) of approximately $2.28$~GB ($2,280,940$~KB) while hosting the 1.6~GB Gemma 2 2B Q4\_K\_M weights and the allocated context buffer. Under these razor-thin margins, expanding $n_{\text{ctx}}$ to 4096 tokens would immediately trigger an out-of-memory SIGKILL from the LMK. 

By enforcing strict JNI-level context budgeting ($n_{\text{ctx}} - 4 - N_{\text{prompt}} - N_{\text{predict}} \ge 0$), HACER-RAG enables mobile applications to operate safely with modest context sizes ($n_{\text{ctx}} \in \{512, 1024\}$) without risking silent truncation or process termination. Furthermore, because HACER-RAG persists KV state to flash and releases the context arena between user interactions, it prevents background RAM bloat, allowing the handset to maintain normal multitasking responsiveness.

\subsection{Thermal and Energy Implications}
Sustained CPU execution at peak clock frequencies (2.2~GHz on Cortex-A78) rapidly elevates device temperature. In continuous evaluation runs, device battery temperature rose from 29.5$^\circ$C to 38.2$^\circ$C, causing the SoC's dynamic voltage and frequency scaling (DVFS) governor to throttle clock rates down to 1.6~GHz.

By eliminating up to \textbf{86.67\%} of prompt-prefill computation ($46.9$~seconds of heavy CPU execution per query in the $K=4$ configuration), HACER-RAG shortens the duration of high-power CPU residency. This dramatic reduction in active processor cycles not only preserves battery life but also prevents thermal throttling, ensuring that subsequent token decoding executes at maximum interactive responsiveness.

\subsection{Architectural Implications for Mobile RAG Systems}
The findings of this paper suggest that on-device RAG systems cannot simply port server-side caching techniques without adapting to mobile realities:
\begin{enumerate}
    \item \textbf{Evidence Binding is Mandatory:} In personal workspaces where local documents are continuously synced or edited, blind session caching is an unacceptable liability. Cryptographic chunk hashing provides an indispensable safety layer with microsecond overhead.
    \item \textbf{Canonical Layout Eliminates Permutation Penalty:} Enforcing canonical evidence ordering overcomes the fragility of causal attention masks, ensuring high cache hit rates across varying retriever score distributions.
    \item \textbf{Selective Downstream Repair Maximizes Efficiency:} By preserving valid prefix KV segments and re-prefilling only modified downstream chunks, runtimes can maintain freshness while preserving the bulk of compute savings.
\end{enumerate}
